% The gullet's macro map is a *cache* of the eqtb; \futurelet (like \let) must
% evict a stale entry when it binds the target to something that is not a macro
% -- a character, here. If it does not, the target keeps looking like a macro,
% and \ifx compares one operand cached-as-macro against a non-macro operand and
% returns false even when their meanings are identical. That is exactly how
% \@ifstar's `\ifx\@let@token\reserved@d' (both \let to `*') came out false,
% so a \DeclareRobustCommand* fed its `*' to \xdef and stalled latex.ltx.
\input plain
\catcode`\@=11
% \tok starts as a macro (so it has a cache entry), then \futurelet rebinds it
% to the character `*'; \ifx against another `*' must then be true.
\def\tok{some macro body}
\let\star=*
\def\peek{\ifx\tok\star \message{<a=YES>}\else \message{<a=NO>}\fi}
\futurelet\tok\peek *
% and the reverse: a macro target \futurelet to a control sequence stays a macro
% comparison (both \relax here) -- true.
\def\tok{x}\let\myrelax\relax
\def\peekb{\ifx\tok\myrelax \message{<b=YES>}\else \message{<b=NO>}\fi}
\futurelet\tok\peekb \relax
\catcode`\@=12
\bye
